\chapter{Blindness}
\index{Blindness}

\section*{Definition and Overview}
Blindness is the complete, or near-complete, loss of sight that cannot be corrected with glasses, contact lenses, or surgery. The term covers a wide spectrum of severity: some people see nothing at all, while others retain only light perception, hand movements, or a narrow tunnel of central vision. Legal blindness is a defined threshold used for benefits and licensing, typically vision of 20/200 or worse in the better eye, or a visual field of 20 degrees or less. Most people who are registered blind retain some useful vision and are more accurately described as having low vision or severe visual impairment.

\section*{Epidemiology}
The World Health Organization estimates that around 40 million people worldwide are blind and that several times that number live with moderate or severe visual impairment. The leading causes differ sharply by region and income: in low-income countries, cataract and trachoma predominate, whereas in high-income countries age-related macular degeneration, glaucoma, and diabetic retinopathy are the main causes. Blindness becomes far more common with age, and overall numbers are rising as populations age. In children, congenital cataract, measles, rubella, and vitamin A deficiency are important preventable causes, and retinopathy of prematurity remains a concern in some settings.

\section*{Aetiology and Pathophysiology}
Blindness results when the eye, the optic nerve, or the visual processing areas of the brain are damaged. The principal mechanisms and causes are:

\begin{itemize}
    \item \textbf{Clouding of the lens:} cataract, which blocks and scatters light and is the commonest cause of blindness worldwide
    \item \textbf{Retinal damage:} age-related macular degeneration and diabetic retinopathy destroy the light-sensitive cells of the retina
    \item \textbf{Optic nerve damage:} glaucoma and other optic neuropathies destroy the nerve that carries visual information to the brain
    \item \textbf{Corneal disease:} trachoma, onchocerciasis, and severe corneal infections or injury
    \item \textbf{Trauma:} penetrating and blunt injuries that damage the eye or visual pathways
    \item \textbf{Brain damage:} stroke, tumours, or injury affecting the occipital lobe and visual cortex
    \item \textbf{Congenital and genetic conditions:} inherited retinal dystrophies, albinism, and developmental abnormalities of the eye
\end{itemize}

\section*{Clinical Features}
Vision loss presents in ways that depend on the site and speed of the damage:

\begin{itemize}
    \item \textbf{Sudden onset:} as in retinal artery occlusion, retinal detachment, optic neuritis, or stroke
    \item \textbf{Gradual onset:} as in cataract, glaucoma, and macular degeneration
    \item \textbf{Central loss:} loss of detailed reading vision, typical of age-related macular degeneration
    \item \textbf{Peripheral loss:} loss of side vision, typical of glaucoma
    \item \textbf{Total or partial, in one or both eyes}
    \item \textbf{Accompanying features:} floaters, flashes of light, eye pain, redness, or distortion of straight lines
\end{itemize}

\section*{Differential Diagnosis}
When a person reports loss of vision, the clinician systematically considers correctable and serious causes:

\begin{itemize}
    \item Refractive error: blurring that is fully correctable with glasses
    \item Cataract: gradual blur and glare from a cloudy lens
    \item Glaucoma: often silent peripheral field loss with raised eye pressure
    \item Age-related macular degeneration: central blur and distortion
    \item Diabetic retinopathy: bleeding, swelling, and new vessels in the retina
    \item Retinal detachment: flashes, floaters, and a curtain across the vision
    \item Optic neuritis: often painful, unilateral vision loss from optic nerve inflammation
    \item Stroke or transient ischaemic attack affecting the visual cortex
    \item Functional (non-organic) vision loss, more common in children and young adults
\end{itemize}

\section*{When to Seek Medical Care}
Any sudden change in vision warrants urgent assessment, even if it resolves completely. Seek immediate review for:

\begin{itemize}
    \item Sudden loss of vision in one or both eyes
    \item Flashes of light or a sudden shower of new floaters
    \item A curtain or shadow spreading across the field of vision
    \item Double vision, or facial weakness and speech problems suggestive of stroke
    \item Eye pain, redness, or light sensitivity accompanied by blurred vision
    \item Gradual vision loss that is beginning to interfere with daily life
\end{itemize}

\textbf{Contact your local emergency services for:} sudden total loss of vision, suspected stroke with visual symptoms, or a serious eye injury.

\section*{Diagnosis and Investigations}
A thorough eye assessment by an optometrist or ophthalmologist is the key to diagnosis:

\begin{itemize}
    \item \textbf{Visual acuity testing:} reading letters on a chart to measure sharpness of sight
    \item \textbf{Visual field testing:} mapping the field of vision to detect peripheral loss
    \item \textbf{Slit-lamp examination:} magnified inspection of the cornea, lens, and front of the eye
    \item \textbf{Fundoscopy:} examination of the retina and optic nerve for diabetic, glaucomatous, and macular changes
    \item \textbf{Tonometry:} measurement of intraocular pressure for glaucoma
    \item \textbf{Imaging:} optical coherence tomography of the retina, and CT or MRI of the brain where stroke or tumour is suspected
    \item \textbf{Electrophysiology:} specialist tests of retinal and optic nerve function in selected cases
\end{itemize}

\section*{Management}
Management depends entirely on the cause, and several forms of blindness are treatable or preventable:

\begin{itemize}
    \item \textbf{Cataract surgery:} replacement of the cloudy lens restores sight in the great majority of cases
    \item \textbf{Glaucoma treatment:} eye drops, laser, or surgery to lower eye pressure and protect the optic nerve
    \item \textbf{Macular degeneration:} anti-VEGF injections can slow and sometimes improve central vision in the wet form
    \item \textbf{Diabetic retinopathy:} laser, intravitreal injections, and tight control of blood sugar, blood pressure, and cholesterol
    \item \textbf{Low-vision aids:} magnifiers, telescopic lenses, high-contrast materials, and improved lighting
    \item \textbf{Rehabilitation:} orientation and mobility training, assistive technology, screen readers, and daily-living skills
    \item \textbf{Support:} vision-loss organisations, peer support, and psychological counselling
\end{itemize}

\section*{Complications}
\begin{itemize}
    \item Falls, fractures, and other injuries from reduced sight
    \item Difficulty managing medicines, cooking, and finances safely
    \item Loss of the ability to drive and reduced independence
    \item Social isolation, anxiety, and depression
    \item Loss of employment and reduced income
    \item Increased demands on carers and family members
\end{itemize}

\section*{Prognosis and Outcomes}
The outlook depends heavily on the underlying cause. Most blindness from cataract is reversible with surgery, and early treatment of glaucoma, diabetic retinopathy, and wet macular degeneration prevents much vision loss. Once vision is lost it is usually permanent, but remaining sight can often be used far more effectively with rehabilitation and low-vision aids. With appropriate support and training, blind and severely visually impaired people can lead full, independent, and productive lives.

\section*{Prevention and Self-Care}
\begin{itemize}
    \item Have regular eye examinations, especially after age 40 and at any age with diabetes
    \item Wear sunglasses to reduce exposure to ultraviolet light
    \item Use protective eyewear for sport, work, and DIY tasks
    \item Control blood sugar, blood pressure, and cholesterol
    \item Stop smoking, which increases the risk of macular degeneration and cataract
    \item Eat a diet rich in fruit and green leafy vegetables
    \item Protect newborns and children from preventable infections and ensure vitamin A adequacy where deficiency is common
    \item Seek prompt care for any sudden change in vision
\end{itemize}

\section*{Sources and Further Reading}
The following sources provide reliable, up-to-date information for patients and students of medicine.

\begin{itemize}
    \item World Health Organization. \textit{Blindness and vision impairment -- facts and figures.} https://www.who.int/news-room/fact-sheets/detail/blindness-and-visual-impairment
    \item NHS. \textit{Overview of vision problems, blindness, and low vision.} https://www.nhs.uk/conditions/vision-loss/
    \item Mayo Clinic. \textit{Blindness and low vision -- causes and care.} https://www.mayoclinic.org/diseases-conditions/blindness/
    \item Centers for Disease Control and Prevention. \textit{Vision health initiative.} https://www.cdc.gov/visionhealth/
    \item MSD Manual. \textit{Vision impairment and blindness -- professional edition.} https://www.msdmanuals.com/
    \item MedlinePlus. \textit{Blindness and vision loss -- patient information.} https://medlineplus.gov/
\end{itemize}
